% Options for packages loaded elsewhere
\PassOptionsToPackage{unicode}{hyperref}
\PassOptionsToPackage{hyphens}{url}
\documentclass[11pt]{article}
\usepackage[a4paper,left=2.7cm,right=2.7cm,top=2.6cm,bottom=2.7cm]{geometry}
\usepackage{xcolor}
\usepackage{amsmath,amssymb}
\setcounter{secnumdepth}{2}
\setcounter{tocdepth}{2}
\usepackage{iftex}
\ifPDFTeX
  \usepackage[T1]{fontenc}
  \usepackage[utf8]{inputenc}
  \usepackage{textcomp} % provide euro and other symbols
\else % if luatex or xetex
  \usepackage{unicode-math} % this also loads fontspec
  \defaultfontfeatures{Scale=MatchLowercase}
  \defaultfontfeatures[\rmfamily]{Ligatures=TeX,Scale=1}
\fi
\usepackage{lmodern}
\ifPDFTeX\else
  % xetex/luatex font selection
\fi
% Use upquote if available, for straight quotes in verbatim environments
\IfFileExists{upquote.sty}{\usepackage{upquote}}{}
\IfFileExists{microtype.sty}{% use microtype if available
  \usepackage[]{microtype}
  \UseMicrotypeSet[protrusion]{basicmath} % disable protrusion for tt fonts
}{}
\makeatletter
\@ifundefined{KOMAClassName}{% if non-KOMA class
  \IfFileExists{parskip.sty}{%
    \usepackage{parskip}
  }{% else
    \setlength{\parindent}{0pt}
    \setlength{\parskip}{6pt plus 2pt minus 1pt}}
}{% if KOMA class
  \KOMAoptions{parskip=half}}
\makeatother
\usepackage{longtable,booktabs,array}
\newcounter{none} % for unnumbered tables
\usepackage{calc} % for calculating minipage widths
% Correct order of tables after \paragraph or \subparagraph
\usepackage{etoolbox}
\makeatletter
\patchcmd\longtable{\par}{\if@noskipsec\mbox{}\fi\par}{}{}
\makeatother
% Allow footnotes in longtable head/foot
\IfFileExists{footnotehyper.sty}{\usepackage{footnotehyper}}{\usepackage{footnote}}
\makesavenoteenv{longtable}
\setlength{\emergencystretch}{3em} % prevent overfull lines
\providecommand{\tightlist}{%
  \setlength{\itemsep}{0pt}\setlength{\parskip}{0pt}}
\usepackage{bookmark}
\IfFileExists{xurl.sty}{\usepackage{xurl}}{} % add URL line breaks if available
\urlstyle{same}
\usepackage{fancyhdr}
\usepackage{enumitem}
\setlist[itemize]{leftmargin=*,itemsep=0.25em,topsep=0.35em}
\setlist[enumerate]{leftmargin=*,itemsep=0.25em,topsep=0.35em}
\definecolor{reportblue}{HTML}{17365D}
\definecolor{reportgray}{HTML}{5B6573}
\pagestyle{fancy}
\fancyhf{}
\fancyhead[L]{\small\textcolor{reportgray}{Critical Review: Five NLP Research Papers}}
\fancyfoot[C]{\thepage}
\setlength{\headheight}{15pt}
\renewcommand{\headrulewidth}{0.3pt}
\renewcommand{\sectionmark}[1]{\markboth{#1}{}}
\hypersetup{
  pdftitle={Critical Review of Five NLP Research Papers},
  pdfsubject={Comparative review of five recent NLP studies},
  hidelinks,
  pdfcreator={LaTeX}}

\begin{document}
\hypersetup{pageanchor=false}
\begin{titlepage}
\thispagestyle{empty}
\centering
\vspace*{2.2cm}
{\color{reportblue}\rule{\textwidth}{1.2pt}\par}
\vspace{1.2cm}
{\Huge\bfseries Critical Review of Five NLP Research Papers\par}
\vspace{0.8cm}
{\Large Objectives, Motivation, Methodologies, Datasets, Evaluation
Metrics, Findings, and Limitations\par}
\vspace{1.2cm}
{\color{reportblue}\rule{0.72\textwidth}{0.6pt}\par}
\vfill
{\large Comparative Review\par}
\vspace{0.35cm}
{\large September 2026\par}
\vspace*{1.8cm}
\end{titlepage}

\pagenumbering{roman}
\hypersetup{pageanchor=true}
\tableofcontents
\clearpage
\pagenumbering{arabic}

\section{Executive Synthesis}\label{executive-synthesis}

Across the five papers, the most consistent result is that carefully
chosen structure and supervision can matter more than raw model scale.
SC-RAG structures evidence before self-correction; 2T-FT concentrates
supervision on two positions; the topic-modeling framework reduces LLM
calls through clustering; the KGC study shows that small LoRA-tuned
models outperform much larger prompting baselines; and PLEM injects
evidence semantics into relation prototypes. The collective lesson is
not that larger models are unnecessary, but that task-specific inductive
design often converts limited compute and data into stronger results.

{\def\LTcaptype{none} % do not increment counter
\begin{longtable}[]{@{}
  >{\raggedright\arraybackslash}p{(\linewidth - 6\tabcolsep) * \real{0.2500}}
  >{\raggedright\arraybackslash}p{(\linewidth - 6\tabcolsep) * \real{0.2500}}
  >{\raggedright\arraybackslash}p{(\linewidth - 6\tabcolsep) * \real{0.2500}}
  >{\raggedright\arraybackslash}p{(\linewidth - 6\tabcolsep) * \real{0.2500}}@{}}
\toprule\noalign{}
\begin{minipage}[b]{\linewidth}\centering
\textbf{Paper}
\end{minipage} & \begin{minipage}[b]{\linewidth}\centering
\textbf{Core intervention}
\end{minipage} & \begin{minipage}[b]{\linewidth}\centering
\textbf{Main evidence}
\end{minipage} & \begin{minipage}[b]{\linewidth}\centering
\textbf{Primary caution}
\end{minipage} \\
\midrule\noalign{}
\endhead
\bottomrule\noalign{}
\endlastfoot
SC-RAG & Hybrid retrieval plus evidence-aware self-correction & Best or
near-best results across varied RAG benchmarks & Single backbone and
task-specific tuning limit generality \\
2T-FT & Loss on first response and first EOS tokens & \textgreater90\%
shorter reported training time than LoRA full-token supervision &
Depends on costly pseudo-label extraction and is evaluated mainly with
greedy decoding \\
Efficient topic modeling & PDDP clusters plus five exemplars per cluster
& About eight minutes across five corpora with strong alignment &
Generator and evaluator share an LLM, creating circularity risk \\
Natural vs programming language KGC & Matched LoRA fine-tuning under
natural and Python-style formats & Roughly 15-21 absolute F1 points over
CodeKGC & No validation inference and strict exact-string scoring
complicate interpretation \\
PLEM & Evidence matching during prototype learning & Average Macro F1
gain of 1.23 points across two benchmarks & Requires evidence
annotations and retains weak cross-domain performance \\
\end{longtable}
}

\subsection{Review Method and Evidence Attribution}\label{review-method-and-evidence-attribution}

The review distinguishes direct evidence from interpretation and
independent critique. This separation prevents reviewer judgments from
being mistaken for claims made by the original authors:

\begin{itemize}
\item
  \textbf{Paper evidence} denotes a point stated directly in the paper,
  including author-acknowledged limitations and future-work directions.
\item
  \textbf{Interpretive synthesis} denotes a close interpretation
  supported by the design, results, or discussion but not presented as a
  formal author claim.
\item
  \textbf{Critical analysis} denotes an independently reasoned
  limitation, risk, or implication and should not be attributed to the
  paper's authors.
\end{itemize}

\clearpage
\section{Towards Evidence-Aware Retrieval-Augmented Generation via
Self-Corrective
Chain-of-Thought}\label{towards-evidence-aware-retrieval-augmented-generation-via-self-corrective-chain-of-thought}

\textbf{Citation.} Li, Y., Ke, W., Liu, J., Wang, P., Liu, J., and He,
Y. Information Processing and Management 63, 104369, 2026. DOI
10.1016/j.ipm.2025.104369.

{\def\LTcaptype{none} % do not increment counter
\begin{longtable}[]{@{}
  >{\raggedright\arraybackslash}p{(\linewidth - 2\tabcolsep) * \real{0.5000}}
  >{\raggedright\arraybackslash}p{(\linewidth - 2\tabcolsep) * \real{0.5000}}@{}}
\toprule\noalign{}
\begin{minipage}[b]{\linewidth}\centering
\textbf{Item}
\end{minipage} & \begin{minipage}[b]{\linewidth}\centering
\textbf{Summary}
\end{minipage} \\
\midrule\noalign{}
\endhead
\bottomrule\noalign{}
\endlastfoot
Objective & Resolve conflicts between retrieved external evidence and
static parametric knowledge without sacrificing inference efficiency. \\
Method & SC-RAG: hybrid semantic/aspect retrieval, LoRA fine-tuning, and
a one-step evidence-aware self-corrective chain of thought. \\
Data & LaMP, HotpotQA, MMLU-Med, LegalBench-RAG, and CodeRAG-Bench. \\
Evaluation & Accuracy, F1, MAE, RMSE, ROUGE-1/L, Pass@1, Precision@k,
Recall@k, time, and significance testing. \\
Headline finding & Reported gains range from 1.0\% to 30.3\%, with
inference-time reductions up to 14.3\%. \\
\end{longtable}
}

\subsection{Objective and Research
Questions}\label{objective-and-research-questions}

\textbf{Paper evidence.} The paper defines its objective in Section 3.3: develop SC-RAG
to reconcile interior-exterior knowledge conflicts while improving
accuracy and factual coherence without extra training data or
inefficient iterative retrieval. The five experimental questions test
overall performance, the separate value of evidence extraction and
self-corrective CoT, the advantage over conventional CoT, sensitivity to
retrieval count and hybrid weight, and whether generated self-correction
exceeds a fixed template.

\subsection{Motivation}\label{motivation}

\textbf{Paper evidence.} Standard RAG can retrieve current information that conflicts
with an LLM\textquotesingle s older parametric knowledge.
Parameter-first methods may require a separate evaluator, while
context-first methods repeatedly retrieve and regenerate, increasing
latency. The authors aim to retain the model\textquotesingle s own
reasoning capacity, obtain finer token-level evidence, and perform a
single evidence-aware correction step (pp. 2-5).

\textbf{Interpretive synthesis.} The work is also motivated by controllability: a system
that exposes the evidence used to revise an answer provides a clearer
operational story than unstructured self-correction, even though the
paper does not establish full explanation faithfulness.

\subsection{Methodology}\label{methodology}

\begin{quote}
\textbf{1.} Retrieve candidate documents with Contriever and compute
semantic relevance.

\textbf{2.} Use an unsupervised attention-based aspect retriever to
identify fine-grained aspect terms and combine the two retrieval signals
using weight alpha.

\textbf{3.} Select supporting documents and evidence tokens, with four
retrieved documents and five aspect terms per sentence in the reported
setup.

\textbf{4.} LoRA-fine-tune Llama-3.1-8B-Instruct to produce
self-corrective reasoning conditioned on the preliminary answer and
evidence.

\textbf{5.} Generate the final response in one correction pass rather
than an iterative retrieve-generate loop.
\end{quote}

\subsection{Datasets}\label{datasets}

{\def\LTcaptype{none} % do not increment counter
\begin{longtable}[]{@{}
  >{\raggedright\arraybackslash}p{(\linewidth - 2\tabcolsep) * \real{0.5000}}
  >{\raggedright\arraybackslash}p{(\linewidth - 2\tabcolsep) * \real{0.5000}}@{}}
\toprule\noalign{}
\begin{minipage}[b]{\linewidth}\centering
\textbf{Dataset}
\end{minipage} & \begin{minipage}[b]{\linewidth}\centering
\textbf{Task and scale stated in the paper}
\end{minipage} \\
\midrule\noalign{}
\endhead
\bottomrule\noalign{}
\endlastfoot
LaMP & Seven personalization tasks; LaMP-6 unavailable and omitted;
individual tasks contain roughly 6K-20K questions. \\
HotpotQA & Multi-hop question answering; approximately 80K questions in
the experimental description. \\
MMLU-Med & Medical multiple-choice questions; 1,089 questions. \\
LegalBench-RAG & Four legal retrieval datasets; 6,858 samples. \\
CodeRAG-Bench & 9,000 code tasks and a 25-million-document retrieval
datastore. \\
\end{longtable}
}

\textbf{Paper evidence.} For these benchmarks, the authors randomly designate 20\% of
training samples as the test set and construct fine-tuning examples with
their self-correction template (Section 5.1).

\subsection{Evaluation Metrics}\label{evaluation-metrics}

\begin{itemize}
\item
  LaMP classification: accuracy, F1, mean absolute error, and root mean
  square error as appropriate to each task.
\item
  LaMP generation: ROUGE-1 and ROUGE-L.
\item
  HotpotQA: answer F1. MMLU-Med: accuracy. CodeRAG-Bench: Pass@1.
\item
  LegalBench-RAG: Precision@k and Recall@k for k = 1, 2, 4, and 8.
\item
  Efficiency: inference time; reliability: five executions per method
  and significance tests reported in the appendix.
\end{itemize}

\subsection{Findings}\label{findings}

\begin{itemize}
\item
  \textbf{Paper evidence.} SC-RAG is best on most reported benchmarks and improves over
  baselines by up to 19.7\% on LaMP-1, 30.3\% on LaMP-2, and 1.7\% on
  LaMP-3 (Tables 1-2).
\item
  \textbf{Paper evidence.} Removing evidence extraction reduces performance by as much
  as 9.9\% on LaMP-2 and 6.2 points on HotpotQA; removing
  self-correction produces still larger reductions (Table 3).
\item
  \textbf{Paper evidence.} Parameter-first baselines are stronger on several QA tasks,
  whereas context-first systems can be stronger for generation; SC-RAG
  attempts to combine their benefits.
\item
  \textbf{Paper evidence.} Four supporting documents provide the best observed
  performance-efficiency balance, and the optimal semantic/aspect
  mixture is task dependent.
\item
  \textbf{Paper evidence.} The method reduces inference time by up to 14.3\% relative
  to selected iterative baselines while avoiding a separate trained
  evaluator.
\end{itemize}

\subsection{Limitations}\label{limitations}

\textbf{Limitations acknowledged by the authors.}

\begin{itemize}
\item
  Specialized medical, legal, and technical language may exceed the
  hybrid retriever\textquotesingle s ability to identify fine-grained
  evidence unless the corpus and retriever are adapted.
\item
  Self-correction can change an initially correct answer into an
  incorrect one; the authors propose evidence thresholds and a fallback
  that retains the original response under ambiguity.
\item
  The hybrid weight alpha is manually selected, and Contriever remains
  frozen while the LLM is tuned. The authors identify adaptive weighting
  and end-to-end retriever training as future work.
\end{itemize}

\textbf{Additional critical observations.}

\begin{itemize}
\item
  Only one 8B backbone is evaluated, so improvements may depend on
  Llama-3.1-8B-Instruct rather than the framework alone.
\item
  The paper selects an optimal retrieval weight for each task but does
  not clearly isolate a validation-only selection protocol;
  task-specific tuning could make test results optimistic.
\item
  Creating test sets from training partitions weakens comparability with
  standard benchmark test protocols and requires especially careful
  leakage controls.
\item
  The self-corrective rationale is treated as useful reasoning, but no
  faithfulness or human evidence-grounding study establishes that the
  textual rationale reflects the model\textquotesingle s actual decision
  process.
\item
  The implementation requires an 80 GB A100 for fine-tuning, which
  qualifies the accessibility claim even though inference is more
  efficient than iterative baselines.
\item
  Table 4 contains a visually anomalous Madam-RAG Precision@8 value of
  0.799 amid values near 0.1; this should be checked against code or
  corrected results before reuse.
\end{itemize}

\subsection{Reproducibility
Assessment}\label{reproducibility-assessment}

Moderate. The backbone, optimizer, learning rate, batch size, retriever,
number of documents, aspect count, and GPU are reported. Reproduction
would still require the complete data-construction code, exact alpha
selection procedure, templates, and baseline versions. Evidence anchors:
Sections 3.3-5.5, Tables 1-4, and Sections 7-8.

\clearpage
\section{2T-FT: Two-Token Fine-Tuning Improves Zero-Shot Performance
With Minimal
Training}\label{t-ft-two-token-fine-tuning-improves-zero-shot-performance-with-minimal-training}

\textbf{Citation.} Neto, P. S., Dyonisio, J. D. S., Lemos, J. F. S. S.,
Kühne, F., Guerra, R. S., and Drews Jr., P. L. J. Neural Computing and
Applications 38, article 321, 2026. DOI 10.1007/s00521-026-12052-9.

{\def\LTcaptype{none} % do not increment counter
\begin{longtable}[]{@{}
  >{\raggedright\arraybackslash}p{(\linewidth - 2\tabcolsep) * \real{0.5000}}
  >{\raggedright\arraybackslash}p{(\linewidth - 2\tabcolsep) * \real{0.5000}}@{}}
\toprule\noalign{}
\begin{minipage}[b]{\linewidth}\centering
\textbf{Item}
\end{minipage} & \begin{minipage}[b]{\linewidth}\centering
\textbf{Summary}
\end{minipage} \\
\midrule\noalign{}
\endhead
\bottomrule\noalign{}
\endlastfoot
Objective & Improve greedy-decoding reasoning by fine-tuning only the
first answer token and first EOS token of implicit CoT paths. \\
Method & CoT-Decoding, BERT answer-span extraction, pseudo-label
selection, and two-position loss masking implemented with LoRA
adapters. \\
Data & GSM8K, MultiArith, SVAMP, plus nine out-of-distribution reasoning
benchmarks. \\
Evaluation & Exact-match accuracy, mean and standard deviation across
three seeds, training time, supervised-token percentage, and task
generalization. \\
Headline finding & Phi-2 gains 22.7 points on MultiArith, 9.0 on GSM8K,
and 2.3 on SVAMP while training time falls by more than 90\%. \\
\end{longtable}
}

\subsection{Objective and Motivation}\label{objective-and-motivation}

\textbf{Paper evidence.} The paper seeks to transfer reasoning paths already present in
top-k decoding into ordinary greedy generation without hand-written CoT
prompts or full-response supervision. The authors hypothesize that the
first response token selects a reasoning trajectory and that the first
EOS token stabilizes termination (Sections 1 and 3.2).

\textbf{Interpretive synthesis.} The method reframes fine-tuning as behavioral routing
rather than knowledge acquisition: it attempts to make the model enter
an existing reasoning basin rather than teach every reasoning token.

\subsection{Methodology}\label{methodology-1}

\begin{quote}
\textbf{1.} For each question, branch on the top k = 10 first-token
alternatives and greedily decode each path.

\textbf{2.} Extract candidate answers with a BERT model rather than
dataset-specific regular expressions.

\textbf{3.} Use self-consistency aggregation and confidence to select an
implicit CoT response as the pseudo-label.

\textbf{4.} Construct a loss mask that supervises only the first
response token and the first EOS or padding token.

\textbf{5.} Train a LoRA adapter and compare it with LoRA full-token
fine-tuning under matched models and datasets.

\textbf{6.} Evaluate the resulting model with greedy decoding, so
inference no longer needs the k-path search.
\end{quote}

\subsection{Datasets and Models}\label{datasets-and-models}

{\def\LTcaptype{none} % do not increment counter
\begin{longtable}[]{@{}
  >{\raggedright\arraybackslash}p{(\linewidth - 2\tabcolsep) * \real{0.5000}}
  >{\raggedright\arraybackslash}p{(\linewidth - 2\tabcolsep) * \real{0.5000}}@{}}
\toprule\noalign{}
\begin{minipage}[b]{\linewidth}\centering
\textbf{Group}
\end{minipage} & \begin{minipage}[b]{\linewidth}\centering
\textbf{Resources}
\end{minipage} \\
\midrule\noalign{}
\endhead
\bottomrule\noalign{}
\endlastfoot
Primary arithmetic & GSM8K, MultiArith, and SVAMP. \\
Generalization & ARC Easy, ARC Challenge, LogiQA, GPQA, PIQA, BBH,
ETHICS, MuTual, and OpenBookQA. \\
Models & Phi-2, Qwen2-1.5B, and Mistral-7B-v0.1 base models. \\
Generation limits & Maximum 300 new tokens; maximum input length 512;
standard QA format. \\
\end{longtable}
}

\subsection{Evaluation Metrics}\label{evaluation-metrics-1}

\begin{itemize}
\item
  Exact-match accuracy under greedy decoding for arithmetic tasks.
\item
  Mean and standard deviation across three random seeds.
\item
  Percentage of supervised response tokens and wall-clock training time.
\item
  Pseudo-label exact-match correctness on 100 sampled instances per
  dataset.
\item
  lm-eval-harness accuracy for the out-of-distribution reasoning suite.
\item
  Ablations over token positions and controlled pseudo-label correctness
  levels.
\end{itemize}

\subsection{Findings}\label{findings-1}

\begin{itemize}
\item
  \textbf{Paper evidence.} On Phi-2, 2T-FT improves GSM8K from 27.7 to 36.7, MultiArith
  from 61.7 to 84.4, and SVAMP from 48.0 to 50.3 exact match (Table 7).
\item
  \textbf{Paper evidence.} On Qwen2-1.5B, MultiArith rises from 43.3 to 75.0; on
  Mistral-7B-v0.1 it rises from 14.3 to 38.33.
\item
  \textbf{Paper evidence.} Mean training time drops from 1 h 15 m to 6 m on Phi-2, 52 m
  to 4 m on Qwen2-1.5B, and 1 h 53 m to 11 m on Mistral (Table 5).
\item
  \textbf{Paper evidence.} Token-position ablation shows that first-token plus EOS
  supervision has the best accuracy-stability trade-off; EOS-only,
  random-token, and answer-span masks perform poorly (Table 12).
\item
  \textbf{Paper evidence.} 2T-FT is comparatively robust to noisy pseudo-labels, but
  gains across non-arithmetic reasoning categories are not uniform.
\end{itemize}

\subsection{Limitations}\label{limitations-1}

\textbf{Limitations acknowledged by the authors.}

\begin{itemize}
\item
  Large training sets can produce overfitting and repetitive first-token
  behavior.
\item
  The method may be poorly suited to code or programmatic solutions,
  where early tokens do not summarize the full reasoning structure.
\item
  CoT-Decoding adds preprocessing cost and pipeline complexity, and
  outcomes depend on extracted-path quality.
\item
  Cross-domain gains are inconsistent, and decoding strategies other
  than greedy decoding are not systematically evaluated.
\end{itemize}

\textbf{Additional critical observations.}

\begin{itemize}
\item
  The \textgreater90\% training-time claim excludes or separates the
  cost of generating k implicit paths and running answer extraction;
  end-to-end cost should include pseudo-label construction.
\item
  The method is tested on three model families and predominantly
  short-answer reasoning tasks, so its effectiveness for long-form
  factual, multilingual, or safety-critical outputs is uncertain.
\item
  A BERT extractor can introduce selection bias: reasoning paths that
  are easy for the extractor to parse may dominate training even when
  other valid forms exist.
\item
  Exact-match improvements do not establish reasoning faithfulness;
  first-token steering may improve answer format or trajectory selection
  without making intermediate reasoning causally valid.
\item
  Comparisons with Rho-1 are limited to one model and three arithmetic
  datasets and do not equalize every training detail.
\end{itemize}

\subsection{Reproducibility
Assessment}\label{reproducibility-assessment-1}

Good. The paper provides code, model names, decoding limits, k, multiple
seeds, detailed result tables, ablations, and configuration appendices.
End-to-end runtime and data-generation storage costs should still be
recorded in a replication. Evidence anchors: Sections 3-5 and Tables
1-17.

\clearpage
\section{Large Language Models for Efficient Topic
Modeling}\label{large-language-models-for-efficient-topic-modeling}

\textbf{Citation.} Theocharopoulos, P. C., Anagnostou, P.,
Georgakopoulos, S. V., Tasoulis, S. K., and Plagianakos, V. P. Neural
Computing and Applications 37, 24421-24439, 2025. DOI
10.1007/s00521-025-11593-9.

{\def\LTcaptype{none} % do not increment counter
\begin{longtable}[]{@{}
  >{\raggedright\arraybackslash}p{(\linewidth - 2\tabcolsep) * \real{0.5000}}
  >{\raggedright\arraybackslash}p{(\linewidth - 2\tabcolsep) * \real{0.5000}}@{}}
\toprule\noalign{}
\begin{minipage}[b]{\linewidth}\centering
\textbf{Item}
\end{minipage} & \begin{minipage}[b]{\linewidth}\centering
\textbf{Summary}
\end{minipage} \\
\midrule\noalign{}
\endhead
\bottomrule\noalign{}
\endlastfoot
Objective & Reduce time and energy for LLM-based topic modeling while
preserving topic quality. \\
Method & Sentence-BERT embeddings, PDDP divisive clustering, five random
representatives per cluster, and Llama-3.1-70B topic generation. \\
Data & BBC, 20 Newsgroups, scientific abstracts, E-commerce reviews, and
U.S. congressional bills. \\
Evaluation & LLM label alignment, TC, TD, AMI, ARI, MIS, harmonic
purity, runtime, energy, and estimated cost. \\
Headline finding & Roughly eight-minute runtimes, higher LLM alignment
than BERTopic, and energy reductions up to 98\% versus document-wise LLM
inference. \\
\end{longtable}
}

\subsection{Objective and Motivation}\label{objective-and-motivation-1}

\textbf{Paper evidence.} The paper aims to discover topics in large unlabeled corpora
without submitting every document to an LLM. Its motivation combines
scalability, operating cost, and environmental impact: cluster cheap
semantic embeddings first, ask the LLM about a few representatives, and
transfer that topic to the cluster (pp. 24421-24424).

\textbf{Interpretive synthesis.} The method treats an LLM as a semantic labeler
positioned after conventional representation and clustering rather than
as the complete topic model. This modularity is the central efficiency
mechanism.

\subsection{Methodology}\label{methodology-2}

\begin{quote}
\textbf{1.} Clean text by removing HTML, whitespace artifacts,
punctuation, numbers, stop words, and by lowercasing.

\textbf{2.} Embed each document into 384 dimensions with Sentence-BERT
all-MiniLM-L6-v2.

\textbf{3.} Apply Principal Direction Divisive Partitioning through
HiPart using PCA, at most 600 splits, and a minimum cluster size of 20.

\textbf{4.} Randomly select five representative texts from each cluster.

\textbf{5.} Prompt locally hosted Llama-3.1-70B-Instruct through vLLM to
return one topic per cluster.

\textbf{6.} Assign the cluster topic to every member and evaluate both
topic alignment and computational efficiency over ten independent
iterations.
\end{quote}

\subsection{Datasets}\label{datasets-1}

{\def\LTcaptype{none} % do not increment counter
\begin{longtable}[]{@{}
  >{\raggedright\arraybackslash}p{(\linewidth - 2\tabcolsep) * \real{0.5000}}
  >{\raggedright\arraybackslash}p{(\linewidth - 2\tabcolsep) * \real{0.5000}}@{}}
\toprule\noalign{}
\begin{minipage}[b]{\linewidth}\centering
\textbf{Dataset}
\end{minipage} & \begin{minipage}[b]{\linewidth}\centering
\textbf{Scale and labels}
\end{minipage} \\
\midrule\noalign{}
\endhead
\bottomrule\noalign{}
\endlastfoot
BBC News & 2,225 articles across five topic areas. \\
20 Newsgroups & 18,846 documents across 20 categories. \\
Scientific abstracts & 20,006 abstracts from computer science,
mathematics, physics, and statistics; sourced from Kaggle. \\
E-commerce reviews & 50,425 texts across four product domains. \\
Bills & 40,661 summaries from the 110th-114th U.S. Congresses with 21
high-level topics. \\
\end{longtable}
}

\subsection{Evaluation Metrics}\label{evaluation-metrics-2}

\begin{itemize}
\item
  LLM-based topic alignment: Llama chooses the best ground-truth label
  for each generated topic.
\item
  Topic quality: Topic Coherence and Topic Diversity.
\item
  Cluster-label agreement: Adjusted Mutual Information, Adjusted Rand
  Index, Mutual Information Score, and Harmonic Purity.
\item
  Efficiency: elapsed seconds, measured kWh, percentage energy
  reduction, and estimated electricity cost in euros.
\item
  Baselines: LLM one-by-one, BERTopic, and TopicGPT on a 7\% sample for
  the latter comparison.
\end{itemize}

\subsection{Findings}\label{findings-2}

\begin{itemize}
\item
  \textbf{Paper evidence.} LLM-alignment scores are 0.861 on BBC, 0.694 on Newsgroups,
  0.820 on E-commerce, 0.784 on Papers, and 0.506 on Bills, all above
  BERTopic (Table 1).
\item
  \textbf{Paper evidence.} The approach finishes in 473-518 seconds across all five
  corpora, whereas one-by-one inference ranges from 2,950 to 74,417
  seconds.
\item
  \textbf{Paper evidence.} The method obtains the strongest overall cluster agreement
  on BBC, E-commerce, and Bills; BERTopic remains stronger on some
  coherence and diversity metrics (Table 2).
\item
  \textbf{Paper evidence.} Reported energy savings reach 96-98\% on several large
  datasets compared with one-by-one inference and 57-75\% against
  BERTopic on four of five datasets.
\item
  \textbf{Paper evidence.} Simplifying 20 Newsgroups labels changes evaluation
  materially, demonstrating that the LLM alignment score is sensitive to
  label wording.
\end{itemize}

\subsection{Limitations}\label{limitations-2}

\textbf{Limitations acknowledged by the authors.}

\begin{itemize}
\item
  Quality and computational cost depend directly on the initial
  clustering and its granularity.
\item
  Coarse clustering may collapse distinct topics; fine clustering
  increases LLM calls and erodes the efficiency advantage.
\item
  A few random representatives may fail to capture outliers, small
  subtopics, or heterogeneous clusters.
\item
  Maximum splits, cluster size, and exemplar count require empirical
  adjustment; automatic selection remains future work.
\end{itemize}

\textbf{Additional critical observations.}

\begin{itemize}
\item
  Llama-3.1-70B generates topics and also maps those topics to reference
  labels. Shared model preferences can inflate alignment and reduce
  evaluator independence.
\item
  Changing Newsgroups labels to broader terms improves scores but also
  changes the target; the revised result should not be treated as
  directly comparable to the original labels.
\item
  The LLM one-by-one baseline is more accurate on every dataset in Table
  1, so the central result is an efficiency-quality trade-off rather
  than unqualified state-of-the-art quality.
\item
  Two A100 80 GB GPUs and 256 GB RAM are substantial resources; the
  framework reduces inference calls but does not by itself establish
  accessibility on ordinary hardware.
\item
  Energy estimates cover measured execution but not model download,
  model loading, infrastructure overhead, or embodied hardware
  emissions.
\item
  Random exemplars introduce variance, yet the framework does not
  compare random selection with medoids, centroids, diversity sampling,
  or active selection.
\end{itemize}

\subsection{Reproducibility
Assessment}\label{reproducibility-assessment-2}

Moderate to good. Models, embeddings, clustering package, major
hyperparameters, prompt, hardware, repeated runs, energy, and cost
assumptions are reported. Strong reproduction would additionally require
released preprocessing code, exact power measurement procedure, random
seeds, and cluster counts per run. Evidence anchors: Sections 3-6,
Tables 1-4, and Figures 1-5.

\clearpage
\section{Natural vs Programming Language in LLM Knowledge Graph
Construction}\label{natural-vs-programming-language-in-llm-knowledge-graph-construction}

\textbf{Citation.} Gajo, P., and Barrón-Cedeño, A. Information
Processing and Management 62(5), 104195, 2025. DOI
10.1016/j.ipm.2025.104195.

{\def\LTcaptype{none} % do not increment counter
\begin{longtable}[]{@{}
  >{\raggedright\arraybackslash}p{(\linewidth - 2\tabcolsep) * \real{0.5000}}
  >{\raggedright\arraybackslash}p{(\linewidth - 2\tabcolsep) * \real{0.5000}}@{}}
\toprule\noalign{}
\begin{minipage}[b]{\linewidth}\centering
\textbf{Item}
\end{minipage} & \begin{minipage}[b]{\linewidth}\centering
\textbf{Summary}
\end{minipage} \\
\midrule\noalign{}
\endhead
\bottomrule\noalign{}
\endlastfoot
Objective & Test whether small open-weight LLMs, model pretraining
domain, prompt representation, rationale training, and model size affect
RDF triple extraction. \\
Method & LoRA-fine-tune 1B-70B natural-language and code-focused models
with natural or Python-style prompts, with or without rationale. \\
Data & CodeKGC subsets of ADE, CoNLL04, and SciERC. \\
Evaluation & Strict complete-triple precision, recall, micro-F1,
class-level analysis, attention analysis, size and LoRA-module
ablations. \\
Headline finding & The best tuned models outperform the strongest
CodeKGC baselines in Table 4 by 17.0 F1 points on ADE, 20.9 on CoNLL04,
and 14.9 on SciERC; prompt language matters little after tuning. \\
\end{longtable}
}

\subsection{Objective and Research
Questions}\label{objective-and-research-questions-1}

\textbf{Paper evidence.} The study asks whether code-pretrained LLMs are better at
structured extraction, whether Python-style prompts outperform
natural-language prompts, whether explicit rationale improves joint
entity and relation extraction, and whether small open-weight models can
exceed closed GPT-3.5 CodeKGC baselines. It also studies model size and
which LoRA modules matter (Sections 1, 4, and 5).

\subsection{Motivation}\label{motivation-1}

\textbf{Paper evidence.} Existing LLM-based KGC research often relies on large
commercial APIs whose weights, training changes, and long-term
availability are outside the researcher\textquotesingle s control. The
authors pursue reproducibility, lower compute, transparency, and reduced
third-party dependence through downloadable open-weight models and
parameter-efficient tuning.

\textbf{Interpretive synthesis.} The natural-versus-code comparison also tests a broader
hypothesis: whether representational syntax supplies a meaningful
structural prior after the model has already observed hundreds of
supervised input-output pairs.

\subsection{Methodology}\label{methodology-3}

\begin{quote}
\textbf{1.} Construct natural-language and Python-style versions of the
same RDF triple-extraction task, with three in-context demonstrations.

\textbf{2.} Optionally require a rationale listing candidate entities
and relations before the final triples.

\textbf{3.} LoRA-tune natural-language models and code models over three
datasets for 200 steps with rank and alpha 16, batch size 8, learning
rate 2e-4, five warm-up steps, and AdamW.

\textbf{4.} Test base instruction models, tuned 7-8B models, 1B and 3B
variants, and 70B variants; the full study covers 175 evaluated model
conditions.

\textbf{5.} Parse predicted triples and compute strict complete-triple
micro metrics averaged over three testing runs.

\textbf{6.} Analyze prompt effects, rationale effects, relation classes,
attention patterns, model size, and LoRA target modules.
\end{quote}

\subsection{Datasets}\label{datasets-2}

{\def\LTcaptype{none} % do not increment counter
\begin{longtable}[]{@{}
  >{\raggedright\arraybackslash}p{(\linewidth - 4\tabcolsep) * \real{0.3333}}
  >{\raggedright\arraybackslash}p{(\linewidth - 4\tabcolsep) * \real{0.3333}}
  >{\raggedright\arraybackslash}p{(\linewidth - 4\tabcolsep) * \real{0.3333}}@{}}
\toprule\noalign{}
\begin{minipage}[b]{\linewidth}\centering
\textbf{Dataset}
\end{minipage} & \begin{minipage}[b]{\linewidth}\centering
\textbf{Train and test documents}
\end{minipage} & \begin{minipage}[b]{\linewidth}\centering
\textbf{Entity and relation structure}
\end{minipage} \\
\midrule\noalign{}
\endhead
\bottomrule\noalign{}
\endlastfoot
ADE & 2,562 / 300 & Drug and disease entities; one adverseEffect
relation in the selected subset. \\
CoNLL04 & 922 / 288 & Three entity types and five relations over news
text. \\
SciERC & 1,366 / 397 & Six entity types, seven positive relations, and a
none class over scientific text. \\
\end{longtable}
}

\textbf{Paper evidence.} The study uses the same subsampled splits as CodeKGC to
preserve baseline comparability. The authors rank SciERC as the most
difficult because of technical language and many relation classes,
CoNLL04 second, and ADE easiest because it has one relation type
(Section 3 and Table 1).

\subsection{Evaluation Metrics}\label{evaluation-metrics-3}

\begin{itemize}
\item
  Strict true positive: head string, relation string, and tail string
  must all match a reference triple exactly.
\item
  Micro precision, recall, and F1 aggregated over complete triples;
  results averaged over three inference runs.
\item
  Per-relation F1 marginalized across models.
\item
  Qualitative error and attention analysis for rationale behavior.
\item
  Ablations for model size and LoRA query, key, value, attention-output,
  and MLP targets.
\end{itemize}

\subsection{Findings}\label{findings-3}

\begin{itemize}
\item
  \textbf{Paper evidence.} The strongest results are split across two
  Mistral variants. Mistral-7B-Instruct-v0.3 reaches 0.812 on ADE with a
  code prompt and 0.705 on CoNLL04 with a natural-language prompt,
  whereas Mistral-7B-v0.3 reaches 0.396 on SciERC with either prompt
  format; all three best conditions omit rationale (Table 4).
\item
  \textbf{Paper evidence.} Relative to the strongest CodeKGC value in
  Table 4 for each dataset, these scores correspond to gains of 17.0,
  20.9, and 14.9 absolute F1 points, respectively. The abstract reports
  17.4, 19.7, and 15.5 points; the difference is retained here as an
  internal reporting discrepancy in the source paper.
\item
  \textbf{Paper evidence.} After fine-tuning, natural and code prompt formats yield
  similar outcomes; before fine-tuning, code prompts can help general
  chat models more noticeably.
\item
  \textbf{Paper evidence.} Rationale training reduces performance in nearly every tuned
  setting. Attention visualizations suggest that entity and relation
  lists disperse attention toward redundant prompt content.
\item
  \textbf{Paper evidence.} Small tuned models outperform 70B models under the same
  1,600-sample and 200-step budget; the authors argue that larger models
  require more data to be steered effectively.
\item
  \textbf{Paper evidence.} Under the fixed 200-step, low-data recipe,
  full LoRA across attention and MLP projections outperforms full
  fine-tuning and restricted-module adapters; updating value projections
  is particularly important (Table 11).
\end{itemize}

\subsection{Limitations}\label{limitations-3}

\textbf{Limitations acknowledged or directly discussed by the authors.}

\begin{itemize}
\item
  The largest models could not be tuned at 16-bit or 8-bit precision and
  were trained with 4-bit QLoRA; the same limited training budget is
  probably insufficient for 70B models.
\item
  DeepSeek Coder\textquotesingle s 4K context restricts experiments with
  many in-context examples.
\item
  The compute required for 168 tuned models and repeated testing leads
  the authors to evaluate directly on the test set without validation
  inference.
\item
  The study does not claim that rationale text is universally harmful;
  its result is specific to the candidate-list rationale design and
  these extraction tasks.
\end{itemize}

\textbf{Additional critical observations.}

\begin{itemize}
\item
  Using the test set without a validation inference stage creates
  selection risk, especially when interpreting the chosen 200-step
  checkpoint, prompt design, and many ablations.
\item
  The natural and code prompts differ in schema length, tokenization,
  system messages, and output grammar, so the comparison is not a
  perfectly isolated test of \textquotesingle language\textquotesingle{}
  alone.
\item
  Strict exact-string matching is reproducible but penalizes equivalent
  aliases and harmless formatting while missing semantic correctness;
  both strict and normalized metrics would be more informative.
\item
  Three small, subsampled datasets and English-only prompts limit
  conclusions about broad knowledge graph construction, long documents,
  multilingual text, and large relation inventories.
\item
  The study tests many conditions but generally omits uncertainty
  intervals from main tables; multiple comparisons increase the chance
  of overinterpreting small differences.
\item
  The article links to public code, but exact reproduction still depends
  on external implementation and environment details that are not fully
  specified in the paper.
\end{itemize}

\subsection{Proposed Replication Configuration}
\label{proposed-replication-configuration}

ADE is the most practical starting point for a controlled replication:
it is the smallest dataset used in the study, contains one relation
type, and has 300 test documents. A resource-conscious protocol uses
Mistral-7B-Instruct-v0.3 with 4-bit LoRA, three demonstrations, matched
natural-language and Python-style outputs, greedy decoding, and strict
micro-F1. Generated Python-style output should be parsed as data and
must never be executed.

\subsection{Reproducibility
Assessment}\label{reproducibility-assessment-3}

Moderate. The paper provides extensive hyperparameters, prompt
appendices, data links, class counts, model names, ablations, and a
public repository. Reproduction is weakened by missing validation
inference, incomplete uncertainty reporting, and dependence on external
implementation details. Evidence anchors: Sections 3-6, Tables 1-11,
and Appendices A-D.

\clearpage
\section{PLEM: Prototype Learning With Evidence Match for Improving
Few-Shot Document-Level Relation
Extraction}\label{plem-prototype-learning-with-evidence-match-for-improving-few-shot-document-level-relation-extraction}

\textbf{Citation.} Pan, D., Sun, Y., Xu, B., Li, J., Yang, Z., Luo, L.,
Lin, H., and Wang, J. IEEE Transactions on Audio Speech and Language
Processing 33, 2505-2514, 2025. DOI 10.1109/TASLPRO.2025.3578781.

{\def\LTcaptype{none} % do not increment counter
\begin{longtable}[]{@{}
  >{\raggedright\arraybackslash}p{(\linewidth - 2\tabcolsep) * \real{0.5000}}
  >{\raggedright\arraybackslash}p{(\linewidth - 2\tabcolsep) * \real{0.5000}}@{}}
\toprule\noalign{}
\begin{minipage}[b]{\linewidth}\centering
\textbf{Item}
\end{minipage} & \begin{minipage}[b]{\linewidth}\centering
\textbf{Summary}
\end{minipage} \\
\midrule\noalign{}
\endhead
\bottomrule\noalign{}
\endlastfoot
Objective & Improve few-shot document-level relation prototypes by
matching entity-pair representations to evidence sentences. \\
Method & Metric meta-learning with task-specific relation, NOTA, MATCH,
and UNMATCH prototypes plus joint relation and evidence losses. \\
Data & FREDo and ReFREDo in 1-DOC and 3-DOC, in-domain and
DocRED-to-SciERC cross-domain settings. \\
Evaluation & Macro F1 over five random seeds, ablations, hyperparameter
sensitivity, support-count analysis, and a limited LLM comparison. \\
Headline finding & Average improvement of 1.17 Macro F1 points on FREDo
and 1.29 on ReFREDo over RAPL, or 1.23 overall. \\
\end{longtable}
}

\subsection{Objective and Motivation}\label{objective-and-motivation-2}

\textbf{Paper evidence.} Few-shot document-level relation extraction must generalize to
unseen relation types from one or three support documents. Earlier
metric methods aggregate whole-document and entity-pair context into
relation prototypes even though only a few sentences provide decisive
evidence. PLEM aims to suppress irrelevant context and adapt evidence
semantics to the relation types in each episode (Introduction and
Section IV).

\textbf{Interpretive synthesis.} The paper views evidence supervision as representation
shaping rather than merely evidence retrieval: the auxiliary task is
valuable because it changes the geometry of relation prototypes used by
the main classifier.

\subsection{Methodology}\label{methodology-4}

\begin{quote}
\textbf{1.} Encode support and query documents with BERT-base and derive
entity-pair contextual representations.

\textbf{2.} Construct relation prototypes from support triples and
task-specific NOTA prototypes for non-relation cases.

\textbf{3.} Create learnable MATCH and UNMATCH evidence prototypes and
adapt them with episode-specific relation semantics.

\textbf{4.} Score sentences as evidence or non-evidence and use
evidence-weighted information to improve the relation representation.

\textbf{5.} Jointly optimize relation classification and evidence
matching in 50,000 sampled meta-learning episodes.

\textbf{6.} Evaluate on unseen relation categories in 1-DOC and 3-DOC,
in-domain and cross-domain episodes.
\end{quote}

\subsection{Datasets}\label{datasets-3}

{\def\LTcaptype{none} % do not increment counter
\begin{longtable}[]{@{}
  >{\raggedright\arraybackslash}p{(\linewidth - 2\tabcolsep) * \real{0.5000}}
  >{\raggedright\arraybackslash}p{(\linewidth - 2\tabcolsep) * \real{0.5000}}@{}}
\toprule\noalign{}
\begin{minipage}[b]{\linewidth}\centering
\textbf{Benchmark}
\end{minipage} & \begin{minipage}[b]{\linewidth}\centering
\textbf{Construction and evaluation}
\end{minipage} \\
\midrule\noalign{}
\endhead
\bottomrule\noalign{}
\endlastfoot
FREDo in-domain & DocRED relation types split into 62 train, 16
development, and 18 test types; 15K test episodes. \\
FREDo cross-domain & Train on DocRED and evaluate 3K episodes over
SciERC. \\
ReFREDo in-domain & Uses Re-DocRED\textquotesingle s expanded 119,991
relational facts with the same relation-type partition; 15K test
episodes. \\
ReFREDo cross-domain & Re-DocRED training with SciERC test episodes,
matching FREDo\textquotesingle s cross-domain setup. \\
\end{longtable}
}

\subsection{Evaluation Metrics}\label{evaluation-metrics-4}

\begin{itemize}
\item
  Macro F1 across relation types, reported as mean and standard
  deviation over five random seeds.
\item
  Early stopping based on development-set Macro F1.
\item
  Ablations for task-specific evidence prototypes, evidence matching,
  and supervised-style evidence retrieval.
\item
  Sensitivity analysis for loss weight, task-specific mixture, and
  number of evidence prototypes.
\item
  Performance stratified by the number of support relation instances.
\item
  Preliminary ChatGPT comparison on 1K sampled episodes per setting
  because of cost and time constraints.
\end{itemize}

\subsection{Findings}\label{findings-4}

\begin{itemize}
\item
  \textbf{Paper evidence.} PLEM exceeds RAPL by an average 1.17 Macro F1 points on
  FREDo and 1.29 on ReFREDo across settings (Table II).
\item
  \textbf{Paper evidence.} Removing evidence matching reduces Macro F1 by 2.33 points
  in-domain and 1.98 cross-domain (Table III).
\item
  \textbf{Paper evidence.} Task-specific evidence prototypes help in-domain performance
  but have almost no effect cross-domain, implying that episode-specific
  evidence semantics do not transfer cleanly.
\item
  \textbf{Paper evidence.} Three support documents outperform one, but gains are not
  monotonic with the number of individual support instances.
\item
  \textbf{Paper evidence.} Supervised DocRE baselines perform poorly in few-shot
  settings because they overfit known relations; the metric-learning
  design better handles unseen categories.
\item
  \textbf{Paper evidence.} GPT-3.5 and GPT-4o improve with three demonstrations but
  remain below metric-based approaches in the sampled comparison; the
  authors explicitly call the comparison not entirely fair.
\end{itemize}

\subsection{Limitations}\label{limitations-4}

\textbf{Attribution note.} The paper has no dedicated limitations
section. The first three constraints below are directly acknowledged in
the experiments and case analysis; the remaining items are independent
critical observations.

\begin{itemize}
\item
  \textbf{Interpretive synthesis.} Cross-domain performance remains substantially below
  in-domain performance, and improved Re-DocRED annotation does not
  resolve domain adaptation.
\item
  \textbf{Interpretive synthesis.} Metric prototypes can be incomplete when a relation
  has only one support instance, causing all tested models to miss
  relevant triples in the case study.
\item
  \textbf{Interpretive synthesis.} Evidence annotations can be missing or incorrect;
  excessive evidence-loss weighting introduces noise and degrades the
  relation prototype.
\item
  \textbf{Additional critical observations.} PLEM depends on evidence sentence annotations in
  the support data, an additional labeling burden that weakens the
  low-resource premise.
\item
  \textbf{Additional critical observations.} Both benchmarks are episode constructions based
  mainly on DocRED/Re-DocRED and SciERC; broader domains, languages, and
  document genres are not tested.
\item
  \textbf{Additional critical observations.} The BERT-base encoder and fixed training recipe
  leave open whether gains persist with stronger document encoders or
  modern long-context models.
\item
  \textbf{Additional critical observations.} Macro F1 is appropriate for relation balance, but
  reporting micro F1, calibration, NOTA error rates, and evidence-match
  quality would better diagnose behavior.
\item
  \textbf{Additional critical observations.} The sampled LLM comparison uses prompt designs
  and inference paradigms that are not equivalent to PLEM and should not
  support a broad conclusion about LLM inferiority.
\end{itemize}

\subsection{Reproducibility
Assessment}\label{reproducibility-assessment-4}

Good on training details: the paper reports encoder, optimizer, learning
rate, warm-up, batch size, clipping, 50K episodes, early stopping,
hyperparameters, GPU, and five seeds. Full reproduction still depends on
the precise benchmark episode sampler, evidence annotations, and
preprocessing pipeline. Evidence anchors: Sections II-VI, Tables I-V,
and Figures 2-4.

\clearpage
\section{Cross-Paper Synthesis and Replication Design}
\label{cross-paper-synthesis-and-replication-design}

\subsection{What the Five Studies Agree
On}\label{what-the-five-studies-agree-on}

\begin{itemize}
\item
  Structure beats indiscriminate scale. Evidence selection, cluster
  representatives, selective token loss, and prototypes all concentrate
  computation on a smaller decision-relevant subset.
\item
  Open or smaller models can be competitive when supervision and task
  formulation are strong, but the most efficient methods often shift
  cost into preprocessing, annotation, or model-specific infrastructure.
\item
  Auxiliary reasoning is not automatically beneficial. The KGC rationale
  hurts most conditions, while SC-RAG benefits from evidence-conditioned
  self-correction; the difference is whether added text supplies
  decision-relevant evidence or redundant tokens.
\item
  Evaluation design is a major source of uncertainty. LLM-as-judge
  circularity, strict string matching, sampled LLM comparisons, custom
  test partitions, and pseudo-label extractors can materially shape
  reported conclusions.
\item
  Cross-domain generalization remains weaker than in-domain performance
  in every paper where it is tested directly or indirectly.
\end{itemize}

\subsection{Recommended Replication}\label{recommended-replication}

The natural-versus-programming-language KGC paper is the best
replication target among the five studies because it has a public data
interface, a clear two-condition comparison, a deterministic evaluation
metric, and a feasible single-dataset experiment. ADE is the recommended
first dataset because it has one relation type, only 300 test documents,
and the strongest paper results, reducing experimental ambiguity before
moving to CoNLL04 or SciERC.

{\footnotesize
\def\LTcaptype{none} % do not increment counter
\begin{longtable}[]{@{}
  >{\raggedright\arraybackslash}p{(\linewidth - 4\tabcolsep) * \real{0.3333}}
  >{\raggedright\arraybackslash}p{(\linewidth - 4\tabcolsep) * \real{0.3333}}
  >{\raggedright\arraybackslash}p{(\linewidth - 4\tabcolsep) * \real{0.3333}}@{}}
\toprule\noalign{}
\begin{minipage}[b]{\linewidth}\centering
\textbf{Decision}
\end{minipage} & \begin{minipage}[b]{\linewidth}\centering
\textbf{Recommended setting}
\end{minipage} & \begin{minipage}[b]{\linewidth}\centering
\textbf{Reason}
\end{minipage} \\
\midrule\noalign{}
\endhead
\bottomrule\noalign{}
\endlastfoot
Dataset & ADE CodeKGC split & Smallest and simplest paper-used
benchmark; 442 test triples in the selected split. \\
Model & Mistral-7B-Instruct-v0.3 & Strongest family in the paper and
feasible with 4-bit LoRA on Kaggle GPU. \\
Conditions & Natural-language list vs Python-style Triple output & Keeps
the semantic task constant while testing representation format. \\
Training & 200 steps, rank/alpha 16, LR 2e-4, three demonstrations &
Matches the central paper configuration while using gradient
accumulation for limited VRAM. \\
Metric & Strict exact complete-triple micro P/R/F1 & Directly comparable
with the paper; add normalized metrics only as secondary analysis. \\
Seeds & At least three & A single Kaggle run is a smoke test, not a
stable replication. \\
\end{longtable}
\normalsize
}

\subsection{Reproduction Package Validation}
\label{reproduction-package-validation}

\begin{itemize}
\item
  Seven unit tests pass for data loading, prompt construction, safe
  parsing, and strict micro-F1.
\item
  An end-to-end deterministic smoke test completed on the ADE test
  split: 300 documents and 442 gold triples were serialized and
  re-parsed correctly in both formats.
\item
  The Kaggle notebook validates inputs before downloading or training a
  model and saves per-example predictions plus metrics for both
  conditions.
\item
  The code-output parser never executes generated Python.
\item
  The completed checks validate data handling and evaluation logic, not
  model quality. A full 7B GPU fine-tuning run remains necessary before
  reporting replicated performance against the paper.
\end{itemize}

\subsection{Scope and Attribution Note}\label{scope-and-attribution-note}

Factual descriptions derive from the five reviewed papers. Statements
identified as interpretive synthesis or critical analysis are reviewer
judgments rather than claims made by the original authors. These labels
should be retained whenever such observations are cited or reused.

\end{document}
